% c08 —— §12 末 + §13 环面镜像构造，p031–p035
% 开头承接 p030 末 (12.23) 的续句「となる。」（不重复上半句）

即上述各式。将这个 $\mathbb{A}^2$ 进一步紧化为 $\PP^2$ 时，新添 3 个轨道
\begin{align}
    \cO_0 &= \{[x_0:x_1:x_2]\in\PP^2 \mid x_0=0,\ x_1\neq 0,\ x_2\neq 0\}, \tag{12.24}\\
    \cO_{01} &= \{[0:0:1]\}, \tag{12.25}\\
    \cO_{02} &= \{[0:1:0]\} \tag{12.26}
\end{align}
它们对应的闭凸锥由
\begin{align}
    \sigma_{\cO_0} &= \R_{\geq 0}\cdot(-1,-1), \tag{12.27}\\
    \sigma_{\cO_{01}} &= \{(n_1,n_2)\in N_\R \mid n_2\leq n_1,\ n_2\leq 0\}, \tag{12.28}\\
    \sigma_{\cO_{02}} &= \{(n_1,n_2)\in N_\R \mid n_2\geq n_1,\ n_1\leq 0\} \tag{12.29}
\end{align}
给出。

另一方面，把 $\mathbb{A}^2$ 或 $\PP^2$ 在原点爆破，则多面锥 $\sigma_{\cO_{12}}$ 被分成 3 个多面锥
\begin{align}
    \{(n_1,n_2)\in\sigma_{\cO_{12}} \mid n_1\geq n_2\}, \tag{12.30}\\
    \{(n_1,n_2)\in\sigma_{\cO_{12}} \mid n_1=n_2\}, \tag{12.31}\\
    \{(n_1,n_2)\in\sigma_{\cO_{12}} \mid n_1\leq n_2\}. \tag{12.32}
\end{align}

对 $N_1$ 上的扇 $\Sigma_1$ 与 $N_2$ 上的扇 $\Sigma_2$，从 $\Sigma_1$ 到 $\Sigma_2$ 的扇的射指的是 Abel 群同态 $f\colon N_1\to N_2$，使得对任意 $\sigma\in\Sigma_1$，存在 $\tau\in\Sigma_2$ 满足 $f(\sigma)\subset\tau$。扇所构成的范畴，与环面簇及（关于稠密环面的同态的）环面同变映射所构成的范畴等价。

环面簇与其上的丰沛环面同变线丛的组 $(X,L)$，称为极化环面簇（polarized toric variety）。对任意极化环面簇 $(X,L)$，存在某个多面体（polyhedron）$\Delta$，使得作为环面的表示有
\begin{equation}
    H^0(L)\cong \C[M\cap\Delta] \tag{12.33}
\end{equation}
成立。这里右端指以 $M\cap\Delta$ 为基的线性空间。若 $X$ 完备，则 $\Delta$ 成为格多面体。反之，对格多面体 $\Delta$，$\bigoplus_{i=0}^{\infty}\C[M\cap i\Delta]$ 自然具有 $\C$ 代数结构，其 $\Proj$ 给出极化环面簇。

多面体与扇处于对偶关系。例如，多面体的余维 1 的面对应于扇的 1 维锥，给出环面簇的环面不变除子。又，多面体的顶点对应于扇的余维 0 的锥，给出环面簇的环面不动点。

\section{环面镜像构造}

设 $N$ 为秩 $d$ 的自由 Abel 群，$M := \Hom(N,\Z)$。对 $M_\R$ 的多面体 $\Delta$，其极对偶（polar dual）由
\begin{equation}
    \Delta^\circ := \{n\in N_\R \mid \text{对任意 } m\in\Delta \text{ 有 } \langle m,n\rangle\geq -1\} \tag{13.1}
\end{equation}
定义。由定义容易看出 $\Delta=(\Delta^\circ)^\circ$。当格多面体 $\Delta$ 的极对偶 $\Delta^\circ$ 再次成为格多面体时，称 $\Delta$ 为自反的（reflexive）[Bat94]。自反多面体的同构类，与 Gorenstein 环面 Fano 簇（即反典范除子为丰沛 Cartier 除子的完备环面簇）的同构类一一对应。

设 $X$ 为 Gorenstein 环面 Fano 簇，环面素除子的集合记为 $I$。把 $I$ 分割为 $I=I_1\coprod\cdots\coprod I_r$，若对任意 $1\leq i\leq r$，$E_i=\sum_{D\in I_i}D$ 都是 nef 的 Cartier 除子，则称这种分割为 nef 分割（nef-partition）[Bor]。一般地，对 Cartier 除子有
{\small
\[
    \text{globally generated}\ \Leftrightarrow\ \text{base point free}\ \Rightarrow\ \text{semi-ample (base point free after some multiple)}\ \Rightarrow\ \text{nef}
\]}
但要注意：在环面簇上，nef 就意味着由整体截面生成。

nef 分割引起 $-K_X$ 的分解
\begin{equation}
    -K_X=E_1+\cdots+E_r \tag{13.2}
\end{equation}
设与 $H^0(\cO_X(E_i))$ 对应的多面体为 $\Delta_i$，则由 (13.2) 得到 $\Delta$ 的 Minkowski 和分解
\begin{equation}
    \Delta=\Delta_1+\cdots+\Delta_r \tag{13.3}
\end{equation}
这称为 $\Delta$ 的 nef 分割。

代数簇 $X$ 上秩为 $r$ 的向量丛 $E$ 的截面 $s$ 在 $x\in X$ 处非退化（non-degenerate），指的是：取 $E$ 的局部标架 $(e_i)_{i=1}^r$，把 $s$ 局部地写成 $\sum_{i=1}^r f_ie_i$ 时，$(f_1,\dots,f_r)$ 构成正则序列。

\begin{theorem}[{[Muk92, Theorem 1.10]}]
光滑代数簇 $X$ 上的向量丛 $E$ 若由整体截面生成，则 $E$ 的一般截面 $s$ 非退化，且 $s$ 的零点 $Y$ 光滑。
\end{theorem}

自反多面体的 nef 分割，作为 $\cO_X(E_1)\oplus\cdots\oplus\cO_X(E_r)$ 的一般截面的零点，给出 $d-k$ 维的 Calabi–Yau 流形。
% [Tutor Note: 原文此处印作 d-k 维，而 nef 分割 (13.3) 有 r 个分量，按上下文或应为 d-r 维；照原文保留。]

固定维数时，由自反多面体的 nef 分割得到的 Calabi–Yau 完全交的 Hodge 数是否有界，这一问题是 [BN08, Question 4.21] 提出的，并由 [BL15] 解决：

\begin{theorem}[{[BL15, Corollary 5.3]}]
对任意 $n$，作为 Gorenstein 环面 Fano 簇的 nef 完全交得到的 $n$ 维 Calabi–Yau 流形双有理有界，从而 Hodge 数也有界。
\end{theorem}

给定自反多面体的 nef 分割 (13.3) 时，对偶 nef 分割（dual nef-partition）由
\begin{equation}
    \nabla_i := \{n\in N_\R \mid \text{对任意 } j \text{ 有 } \langle \Delta_j,n\rangle\geq -\delta_{ij}\} \tag{13.4}
\end{equation}
定义，并且
\begin{equation}
    \Delta^\circ=\operatorname{Conv}(\nabla_1\cup\cdots\cup\nabla_r) \tag{13.5}
\end{equation}
成立。进一步，
\begin{equation}
    \nabla:=\nabla_1+\cdots+\nabla_r \tag{13.6}
\end{equation}
是自反多面体的 nef 分割，且
\begin{equation}
    \nabla^\circ=\operatorname{Conv}(\Delta_1\cup\cdots\cup\Delta_r) \tag{13.7}
\end{equation}
成立。由 nef 分割定义的 Calabi–Yau 流形与由对偶 nef 分割定义的 Calabi–Yau 流形被期待互为镜像。这样构造出的镜像（的候选），在 $r=1$（超曲面）的情形称为 Batyrev 镜像，对一般的 $r$ 称为 Batyrev–Borisov 镜像。

与自反多面体相伴的环面 Fano 簇一般具有 Gorenstein 奇点。通过在不改变 1 维锥的前提下把扇细分到成为单纯扇为止，可以在不改变典范除子的情况下使其成为只具有 Gorenstein 商奇点的环面弱 Fano 簇，但在 4 维及以上的情形，Gorenstein 商奇点一般不具有 crepant 奇点消解。这种环面簇中的 Calabi–Yau 完全交，具有源自作为容器的环面簇的奇点的商奇点，而 Vafa 公式在这种情况下也给出应当成为 Hodge 数的数。

关于把 Vafa 公式所给出的「Hodge 数」在数学上定式化、并证明拓扑镜像对称这一问题，最著名的结果大概是 [BB96, Theorem 4.15]。其中证明了：由自反多面体的 nef 分割得到的 Batyrev–Borisov 镜像对 $(Y,\check Y)$，对于 [BD96] 对至多具有 Gorenstein toroidal 奇点的流形所引入的弦论 Hodge 数（string-theoretic Hodge number），满足
\begin{equation}
    h^{p,q}_{\mathrm{st}}(Y)=h^{p,n-q}_{\mathrm{st}}(\check Y) \tag{13.8}
\end{equation}
然而，弦论 Hodge 数一般认为并不是「正确」的概念：

% [Tutor Note: 原文 §13 的「定理13.1, 13.2」与「問題13.3, 定義13.4, 定義13.5, 例13.6, 例13.7」共用同一编号序列；
%  preamble 已按此把定理/定义/例/注/问题定义为共用一个计数器，编号自动与原文一致。]

\begin{problem}[{[BM03, Remark 2.9]}]
试给出一个具有交换稳定化子的轨形，使它的轨形 Hodge 数与粗模概形的弦论 Hodge 数不一致。
\end{problem}

作为弦论 Hodge 数的改良，弦 Hodge 数（stringy Hodge number）在 [Bat98] 中对至多具有对数终端奇点的代数簇定义。

环面簇的子簇若与环面轨道横截相交，则称为非退化（non-degenerate）的。对 Gorenstein 环面簇的非退化完全交子簇，弦论 Hodge 数与弦 Hodge 数一致 [BM03, Corollary 2.7]。

弦论 Hodge 数与弦 Hodge 数不同的例子，可在 [BM03, Remark 2.8] 中看到。

通过把 Kontsevich 的母题积分理论 [DL99, Cra04] 推广到轨形，已知轨形的轨形 Hodge 数与粗模空间的弦 Hodge 数一致 [Yas04]。

作为包含 Batyrev–Borisov 镜像以及后文将述的刚性 Calabi–Yau 流形的镜像的框架，[BB97] 引入了下述概念：

\begin{definition}[{[BB97]}]
对 $M_\R$ 的 $d$ 维有理强凸多面锥 $K$，若存在某个 $\deg^\vee\in N$，使得 $K$ 的所有原始生成元都位于由 $\langle -, \deg^\vee\rangle=1$ 定义的 $M_\R$ 的超平面上，则称 $K$ 为 Gorenstein 锥（Gorenstein cone）。$K$ 的对偶锥 $K^\vee:=\{m\in M_\R \mid \text{对任意 } n\in N \text{ 有 } \langle m,n\rangle\geq 0\}$ 再次成为 Gorenstein 锥时，称 $K$ 为自反 Gorenstein 锥（reflexive Gorenstein cone）。对自反 Gorenstein 锥 $K$，称 $K$ 的次数元 $\deg^\vee\in N$ 与 $K^\vee$ 的次数元 $\deg\in M$ 的配对 $\langle \deg^\vee,\deg\rangle\in\Z$ 为 $K$ 的指数（index）。
% [Tutor Note: 原文此处用「対合」（对合，involution）一词，按文意为「対」（配对），译文按文意作「配对」；对照 (13.9) 的同一配对可佐证。]
% [Tutor Note: 原文 K 的双对锥写作 \{m\in M_R \mid 任意の n\in N に対して …\}，但 m\in M_R 与「K∨ 的次数元 deg\in M」不相容（K∨ 应在 N_R 中，此处应为 n\in N_R 且对任意 m\in M）；照原文保留，请审校确认。]
\end{definition}

\begin{exercise}
证明：Gorenstein 锥的次数元唯一确定。
\end{exercise}

对自反 Gorenstein 锥 $K\subset M_\R$，
\begin{equation}
    \Delta:=\{m\in K \mid \langle m,\deg^\vee\rangle=1\},\qquad \Delta^\vee:=\{n\in K^\vee \mid \langle \deg,n\rangle=1\} \tag{13.9}
\end{equation}
成为格多面体，把其格点的集合记为
\begin{equation}
    K_{(1)}:=\Delta\cap M,\qquad K^\vee_{(1)}:=\Delta^\vee\cap N \tag{13.10}
\end{equation}

设 $\Delta\subset M_\R$ 为格多面体，$r$ 为正整数。当 $r\Delta$ 的内部含有格点 $m$，且 $r\Delta-m$ 是自反的时，称 $\Delta$ 为指数 $r$ 的 Gorenstein 多面体（Gorenstein polytope of index $r$）。

\begin{exercise}[{[BB97, Proposition 2.11]}]
设 $K\subset M_\R$ 为 Gorenstein 锥。证明：$K$ 为指数 $r$ 的自反 Gorenstein 锥，当且仅当
\begin{equation}
    \Delta:=\{m\in M_\R \mid \langle m,\deg\rangle=1\} \tag{13.11}
\end{equation}
为指数 $r$ 的 Gorenstein 多面体。
% [Tutor Note: 原文 (13.11) 写作 \langle m,\deg\rangle，而 deg\in M、m\in M_R，二者无法配对；对照 (13.9)，此处或应为 \deg^\vee；照原文保留，请审校确认。]
\end{exercise}

对 $\overline{M}_\R$ 的多面体 $\Delta_1,\dots,\Delta_r$，其 Cayley 多面体（Cayley polytope）由
\begin{equation}
    \Delta_1*\cdots*\Delta_r:=\operatorname{Conv}\{\Delta_1\times e_1,\cdots,\Delta_r\times e_r\} \tag{13.12}
\end{equation}
定义。这里 $M_\R:=\overline{M}_\R\oplus\R^r$，$e_i$ 是 $\R^r\subset M_\R$ 的第 $i$ 个标准基。由 Cayley 多面体生成的 $M_\R$ 的锥称为 Cayley 锥（Cayley cone）。

$\Delta$ 的两个 nef 分割 $\Delta=\Delta_1+\cdots+\Delta_r$ 与 $\Delta=\Delta'_1+\cdots+\Delta'_r$ 若通过 $M$ 的元素列 $m_1,\dots,m_r$ 以 $\Delta'_i:=\Delta_i+m_i$ 相联系，则它们所定义的一般 Calabi–Yau 完全交同构，但对应的对偶 nef 分割
\begin{equation}
    \nabla=\nabla_1+\cdots+\nabla_r \tag{13.13}
\end{equation}
与
\begin{equation}
    \nabla'=\nabla'_1+\cdots+\nabla'_r \tag{13.14}
\end{equation}
所伴随的 Calabi–Yau 完全交一般不同构。这称为双镜像（double mirror）。

下述概念为 Batyrev–Borisov 镜像提供了比 nef 分割「更灵活且概念上更优」（[BL18, Section 2.2] 开头的话）的视角。

\begin{definition}[{[BN08]}]
对指数 $r$ 的自反 Gorenstein 锥 $K\subset M_\R$，若 $K$ 的次数元 $\deg^\vee$ 可由适当的 $t_1,\dots,t_r\in K^\vee_{(1)}$ 表示为
\begin{equation}
    \deg^\vee=t_1+\cdots+t_r \tag{13.15}
\end{equation}
则称 $K$ 完全分解（completely split），并称 (13.15) 为 $\deg^\vee$ 的分解（decomposition of $\deg^\vee$）[BN08]。
\end{definition}

$\deg^\vee$ 的分解引起 $K_{(1)}$ 的分解
\begin{equation}
    K_{(1)}=\coprod_{i=1}^{k}K_{(1),i},\qquad K_{(1),i}:=\{m\in K_{(1)} \mid \langle x,t_j\rangle=\delta_{ij}\} \tag{13.16}
\end{equation}
$K_{(1),i}$ 可以看作 $\bigcap_{i=1}^{t}\operatorname{Ann}t_i$ 的子集，而非 $M$ 的子集。当 $K$ 是 $r$ 个多面体 $\Delta_1,\dots,\Delta_r\subset\overline{M}_\R$ 的 Cayley 锥时，$\bigcap_{i=1}^{t}\operatorname{Ann}t_i$ 与 $\overline{M}$ 自然等同，在此等同之下 $K_{(1),i}$ 与 $\Delta_i\cap\overline{M}$ 一致。
% [Tutor Note: (13.16) 中分解式的指标上限印作 k，而 (13.15) 的分解有 r 个分量；又 \langle x,t_j\rangle 中的 x 在此前未被定义，按同一式应为 m；另，正文两处 \bigcap_{i=1}^{t} Ann t_i 的上限字母 t 与指标 i 混用，按 (13.15) 上限或应为 r。均照原文保留，请审校确认。]

\begin{example}[{[BB97, Example 3.11]}]
设 $\R^d$ 的第一象限为 $\sigma$，并假定 $d=kr$。令 $\Z^d$ 的部分 Abel 群 $M$ 为
\begin{equation}
    M:=\{(m_1,\dots,m_d)\in\Z^d \mid m_1+\cdots+m_d\equiv 0 \mod k\} \tag{13.17}
\end{equation}
则 $\sigma$ 成为 $M$ 的指数 $r$ 的 Gorenstein 锥，但不能表成 2 个格多面体 $\Delta_1,\Delta_2\subset\overline{M}_\R$ 的 Cayley 和
\begin{equation}
    \{(\lambda m_1+\mu m_2,\lambda,\mu)\in M_\R:=\overline{M}_\R\times\R^2 \mid m_1\in\Delta_1,\ m_2\in\Delta_2,\ \lambda,\mu\in\R^{\geq 0}\} \tag{13.18}
\end{equation}
的形式。事实上，若能表成 (13.18) 的形式，则对 $\Delta_1$ 的顶点 $v_1$ 与 $\Delta_2$ 的顶点 $v_2$，联结 $m_1:=(v_1,1,0)$ 与 $m_2:=(v_2,0,1)$ 的线段内部不含格点。另一方面，对 $\sigma$ 的任意生成元 $m_1,m_2$，联结 $m_1$ 与 $m_2$ 的线段内部含有 $k-1$ 个 $M$ 格点。
\end{example}

\begin{example}[{[BB97, Proposition 5.5]}]
设 $K\subset\R^{3r}$ 为第一象限，令
\begin{equation}
    M:=\Z^{3d}+\Z\cdot\left(\frac13,\dots,\frac13\right) \tag{13.19}
\end{equation}
则 $K$ 成为指数 $r$ 的自反 Gorenstein 锥。$\deg^\vee=(1,\dots,1)$，令 $t_i=e_{3i}+e_{3i+1}+e_{3i+2}$，则
\begin{equation}
    \deg^\vee=t_1+\cdots+t_r \tag{13.20}
\end{equation}
给出 $\deg^\vee$ 的分解。与例 13.6 合起来可知：即使自反 Gorenstein 锥完全分解，其对偶锥也未必完全分解。
% [Tutor Note: (13.19) 印作 \Z^{3d}，而 K\subset\R^{3r}、t_i 用到 3r 个坐标，按上下文或应为 \Z^{3r}；照原文保留，请审校确认。]
\end{example}

[BN08, Question 5.2] 中提出了由 $\deg^\vee$ 的两种分解得到的双镜像是否双有理等价的问题，[Li16] 给出了肯定解决。又，[BN08, Conjecture 5.3] 中猜测由 $\deg^\vee$ 的两种分解得到的双镜像导出等价，[FK17] 中给出了肯定解决。

[BL18, Definition 3.1] 中提倡把下面定义的 $D(K,c)$ 看作与 Gorenstein 锥相伴的范畴：选取并固定一个以 $K^\vee$ 为支撑的正则（regular/coherent）且单纯（simplicial）的扇 $\Sigma$，使得其所有 1 维锥的原始生成元都含于 $K^\vee_{(1)}$。把由 $\Sigma$ 所决定的 $\C^{K^\vee_{(1)}}$ 的 Stanley–Reissner 轨迹的补集记为 $U_\Sigma$。$\C^{K^\vee_{(1)}}$ 的点 $z\colon K^\vee_{(1)}\to\C$ 属于 $U_\Sigma$ 的充要条件是：存在 $\Sigma$ 的锥 $\sigma$，使得 $z$ 的零点集是 $\sigma$ 的子集。环面
\begin{equation}
    L:=\left\{\lambda\in(\mathbb{G}_m)^{K^\vee_{(1)}} \;\middle|\; \text{对任意 } m\in M \text{ 有 } \prod_{n\in K^\vee_{(1)}}\lambda(n)^{\langle m,n\rangle}=1\right\} \tag{13.21}
\end{equation}
自然作用于 $U_\Sigma$；关于这个作用的商栈记为 $\cX_\Sigma:=[U_\Sigma/L]$。对 $c\in\C^{K_{(1)}}$，令
\begin{equation}
    C(z):=\sum_{m\in K_{(1)}}c(m)\prod_{n\in K^\vee_{(1)}}z(n)^{\langle m,n\rangle} \tag{13.22}
\end{equation}
进一步，把包含 $G$ 作为子群的更大群 $\widehat G$ 定义为
\begin{equation}
    \widehat G:=\left\{\lambda\in(\mathbb{G}_m)^{K^\vee_{(1)}} \;\middle|\; \text{对任意 } m\in\operatorname{Ann}(\deg^\vee) \text{ 有 } \prod_{n\in K^\vee_{(1)}}\lambda(n)^{\langle m,n\rangle}=1\right\} \tag{13.23}
\end{equation}
固定一个满足 $\langle m,\deg^\vee\rangle=1$ 的 $m\in M$，则映射
\begin{equation}
    \widehat G\to\mathbb{G}_m,\qquad \lambda\mapsto \prod_{n\in K^\vee_{(1)}}\lambda(n)^{\langle m,n\rangle} \tag{13.24}
\end{equation}
是满射，其核不依赖于 $m$ 的取法而与 $G$ 一致。在此情形下，把 $[U_\Sigma/\widehat G]$ 上函数 $C$ 的矩阵因子化所构成的范畴记为 $D(K,c)$，则它只依赖于 $K$ 与 $c$，不依赖于 $\Sigma$ 的取法。当 $K$ 来自 Gorenstein 多面体的 nef 分割时，$D(K,c)$ 与对应的 Calabi–Yau 完全交的导出范畴等价。
% [Tutor Note: 原文 (13.21) 定义的环面记作 L，而其后两处「G を部分群として含む」「…G と一致する」中的 G 未另行定义，按上下文（核应为 L）或即 L；照原文保留，请审校确认。]

% CHUNK c08 DONE pages p031–p035
